\chapter{Quick start: from a diagram to a delivered configuration}\label{ch:quickstart}

This chapter walks through a first delivery end to end with the example organisation that ships
with the kit: a \emph{payments} business unit with a \emph{checkout} application in
\texttt{prod} and \texttt{nonprod}, human access through Entra ID (OIDC) and workloads on AWS and
Kubernetes. Replace the names with yours as you go.

\section{Step 1 -- Describe the organisation to your AI agent}
\begin{session}{Design from a description}
\you \texttt{akeyless-watchover design} Payments BU, apps checkout and refunds, environments prod
and nonprod. Payments admins (Entra group grp-akl-payments-admins) manage nonprod and read prod.
Checkout prod pods (EKS namespace payments, service account payments-api) read their own prod
secrets. Auditors list everything. Break-glass secrets are under prod/break-glass and nobody
but the platform team reads them.

\agent Returns five artefacts: (1) the path tree
\texttt{/payments/<app>/<env>/...}; (2) an access matrix persona~$\times$~path; (3) the YAML model
(\texttt{model/roles/payments.yaml}, \texttt{model/auth\_methods.yaml}, \texttt{model/folders.yaml});
(4) \texttt{access\_expectations.yaml} entries that prove the matrix; (5) what the plan must show,
and the equivalent CLI commands.
\end{session}

\section{Step 2 -- Put the model in the repository}
A role file is short. Rules are path patterns with capabilities; associations bind the role to an
auth method, narrowed by sub-claims:
\begin{lstlisting}
roles:
  /governance/roles/payments/team-admin:
    description: Payments engineers manage nonprod secrets of all their apps, read prod
    audit_access: scoped
    rules:
      - { path: "/payments/+/nonprod/*", capabilities: [read, list, create, update, delete] }
      - { path: "/payments/+/prod/*", capabilities: [read, list] }
      - { path: "/payments/+/prod/break-glass/*", capabilities: [deny] }
    associations:
      - auth_method: /governance/auth/sso-entra-oidc
        sub_claims:
          groups: [grp-akl-payments-admins]
\end{lstlisting}
Every key is documented in Chapter~\ref{doc-model-reference}. The expectations file states, per
persona, the exact capabilities on sample paths (\texttt{[]} means no access):
\begin{lstlisting}
checks:
  - { persona: payments-admin, path: /payments/checkout/prod/db/app,           expect: [read, list] }
  - { persona: payments-admin, path: /payments/checkout/prod/break-glass/root, expect: [] }
\end{lstlisting}

\section{Step 3 -- Test offline}
\begin{lstlisting}
scripts/watchover simulate --model terraform/model \
    --expectations terraform/model/access_expectations.yaml --explain
scripts/watchover test --tool tofu
\end{lstlisting}
\texttt{simulate} evaluates the expectations against the model (wildcards, deny precedence,
sub-claims); \texttt{test} runs everything CI runs, including the Terraform guards that reject
unsafe models (a root-scope rule, a deny mixed with other capabilities, an association without
sub-claims to an auth method that forces them, \ldots). Nothing touches Akeyless.

\section{Step 4 -- Bootstrap the CI identity (once)}
An Akeyless administrator creates the identity the pipeline runs as. It is keyless: the CI's own
OIDC token is trusted, scoped to your repository (plan) and to its protected main branch (apply).
\begin{lstlisting}
cd bootstrap
cp bootstrap.tfvars.example bootstrap.tfvars          # platform, repository, gitlab_url...
export AKEYLESS_ACCESS_ID=... AKEYLESS_ACCESS_KEY=... # admin API key, used once
terraform init && terraform apply -var-file=bootstrap.tfvars
terraform output pipeline_access_id                    # -> CI variable AKEYLESS_PIPELINE_ACCESS_ID
\end{lstlisting}
Then create the CI variables: \texttt{AKEYLESS\_PIPELINE\_ACCESS\_ID}, and as masked JSON maps
\texttt{TF\_VAR\_auth\_method\_secrets}, \texttt{TF\_VAR\_event\_forwarder\_secrets},
\texttt{TF\_VAR\_gateway\_secrets}. Optional: \texttt{AKEYLESS\_API}, \texttt{AKEYLESS\_GATEWAY\_API}.

\section{Step 5 -- Merge request, plan, apply}
\begin{enumerate}
\item Push a branch and open a merge request. CI runs the checks and a \texttt{plan} with the plan
  identity. CODEOWNERS asks the business unit and the platform team to approve.
\item Read the plan. Creations and in-place updates are expected; \emph{replacements} of roles or
  auth methods are red flags (names are ForceNew; an auth method replacement changes its access ID).
  Ask your agent: \texttt{akeyless-watchover review <paste the plan>}.
\item Merge. The apply job runs behind a protected environment, then applies the account settings
  the provider cannot manage (\texttt{watchover account-settings}).
\end{enumerate}

\section{Step 6 -- Prove the result live}
Give each test persona a CLI profile (\texttt{akeyless configure --profile <name> ...}), add
\texttt{profile: <name>} to it in \texttt{access\_expectations.yaml}, and run:
\begin{lstlisting}
scripts/watchover verify --expectations terraform/model/access_expectations.yaml
\end{lstlisting}
The nightly drift job then keeps the account equal to Git: a console change makes the plan
non-empty and the job fails; the next apply reverts it.

\section{Step 7 -- Add notifications}
\begin{lstlisting}
scripts/watchover forwarder --name /governance/events/secops-mail --type email \
    --email-to secops@example.com --preset critical --items-source '/*' --auth-method-source '/*' \
    --write terraform/model/governance.yaml
\end{lstlisting}
Commit, merge request, merge: the forwarder is delivered like everything else.

\begin{note}[Already have an Akeyless account full of configuration?]
Start with \texttt{watchover export} instead of Step~2 (Chapter~\ref{doc-brownfield-import}):
it writes the model and the import blocks, and the first plan should show imports only.
\end{note}
